\documentclass[10.5pt,a4paper]{article}
\usepackage[utf8]{inputenc}
\usepackage[T1]{fontenc}
\usepackage{polski}
\usepackage[top=1.8cm,bottom=1.8cm,left=2cm,right=2cm]{geometry}
\usepackage{hyperref}
\usepackage{xcolor}
\usepackage{booktabs}
\usepackage{tabularx}
\usepackage{listings}
\usepackage{enumitem}
\usepackage{titlesec}

\definecolor{primary}{RGB}{24, 76, 120}
\definecolor{secondary}{RGB}{41, 128, 185}
\definecolor{accent}{RGB}{39, 174, 96}

\hypersetup{
    colorlinks=true,
    linkcolor=primary,
    urlcolor=secondary
}

\titleformat{\section}{\color{primary}\normalfont\large\bfseries}{\thesection.}{0.8em}{}[\color{primary}\titlerule]
\titleformat{\subsection}{\color{secondary}\normalfont\normalsize\bfseries}{\thesubsection.}{0.8em}{}
\titlespacing*{\section}{0pt}{1.2ex plus 0.5ex minus 0.2ex}{0.8ex plus 0.2ex}
\titlespacing*{\subsection}{0pt}{1.0ex plus 0.3ex minus 0.2ex}{0.5ex plus 0.2ex}

\lstset{
    basicstyle=\ttfamily\footnotesize,
    backgroundcolor=\color{black!4},
    frame=single,
    rulecolor=\color{black!15},
    breaklines=true,
    tabsize=2,
    showstringspaces=false
}

\begin{document}

\begin{center}
    {\Large \textbf{\color{primary}System Obsługi Sieci Paczkomatów}}\\[0.15cm]
    {\large \textbf{Dokumentacja Rozbudowy o Moduł Zwrotów Częściowych}}\\[0.1cm]
    {\footnotesize Inżynieria Oprogramowania --- Etap 3 (Wymiana Projektów) --- Grupa ,,Krzycha''}\\[0.25cm]
    \small \textbf{Autorzy:} Filip Kontek, Marcel Kosmala, Igor Krzystek \quad$\vert$\quad \textbf{Prowadzący:} dr inż. Marek Zachara\\
    \footnotesize Akademia Górniczo-Hutnicza w Krakowie, Wydział Elektrotechniki, Automatyki, Informatyki i Inżynierii Biomedycznej
\end{center}

\vspace{-0.1cm}
\hrule
\vspace{0.25cm}

\section{Wstęp i Zakres Prac}

W ramach trzeciego etapu projektu laboratoryjnego zespół otrzymał do analizy, przetestowania i rozbudowy system obsługi paczkomatów autorstwa grupy ,,Krzycha''. Zadanie zostało podzielone na dwa ściśle określone obszary:
\begin{enumerate}[leftmargin=*,itemsep=1pt]
    \item \textbf{Zadanie 1 (Weryfikacja procesu nadania i dostarczenia):} Opracowanie zestawu scenariuszy testowych weryfikujących pełny cykl życia przesyłki z uwzględnieniem sytuacji awaryjnych (anulowanie na etapie zamówienia/nadania, pusta skrytka przy odbiorze przez kuriera, timeouty tranzytowe, błędy odbioru). Każdy test zorganizowano w czytelnym podziale na trzy sekcje: \textit{Arrange} (przygotowanie), \textit{Act} (wywołanie metod biznesowych) oraz \textit{Assert} (weryfikacja stanu).
    \item \textbf{Zadanie 2 (Nowa funkcjonalność --- Obsługa zwrotów w tym częściowych):} Rozbudowa systemu o możliwość wyznaczenia zawartości przesyłki oraz zgłoszenia i realizacji zwrotu jednego lub kilku zamówionych produktów w ramach odebranej przesyłki.
\end{enumerate}

\section{Decyzje Architektoniczne i Eliminacja Redundancji}

W toku analizy inżynierskiej zrezygnowano ze sztucznego podejścia opartego o osobną tabelę ze zdublowaną logistyką tranzytu (antywzorzec spotykany w generatorach kodu). Wdrożono rozwiązanie oparte o czyste zasady \textbf{Domain-Driven Design (DDD)} oraz zasadę \textbf{DRY (Don't Repeat Yourself)}:

\begin{itemize}[leftmargin=*,itemsep=2pt]
    \item \textbf{Rozróżnienie jednostki logistycznej od handlowej:}
    W sieci paczkomatów przesyłka (\texttt{Package}) to jednostka transportowa (gabaryt, rezerwacja skrytek, transport kurierski, kod odbioru). Zawartość paczki to niezależne produkty (\texttt{PackageItem}). Klient nie zwraca części fizycznego kartonu, lecz konkretne pozycje ze środka.
    \item \textbf{Zwrot jako nowa przesyłka powrotna:}
    Fizyczny proces zwrotu podlega dokładnie tym samym prawom logistycznym co pierwotne nadanie. Zamiast tworzyć osobny serwis kurierski dla zwrotów, serwis zwrotów (\texttt{PackageReturnService}) tworzy nową instancję \texttt{Package} za pośrednictwem istniejącego \texttt{PackageOrderService}, automatycznie alokując wolne skrytki i odwracając paczkomaty (nadanie w paczkomacie odbiorcy, dostawa do paczkomatu sklepu).
    \item \textbf{Płynna obsługa zwrotów częściowych i sekwencyjnych:}
    Każdy produkt (\texttt{PackageItem}) posiada powiązanie \texttt{ReturnPackageId}. Jeśli klient zwraca jeden produkt, zostaje on przypisany do nowej paczki zwrotnej, a pozostałe produkty pozostają u klienta. Klient może po tygodniu zgłosić kolejny zwrot drugiego produktu --- system utworzy dla niego drugą, niezależną przesyłkę zwrotną.
    \item \textbf{Brak dualizmu stanów:}
    Oryginalna paczka zachowuje historyczny status \texttt{PackageStatus.Collected} (nie zakłócając pierwotnego cyklu życia), udostępniając właściwości pomocnicze \texttt{HasReturns} oraz \texttt{AreAllItemsReturned}.
\end{itemize}

\section{Wprowadzone Zmiany w Strukturze Kodu}

\begin{table}[h!]
\centering
\footnotesize
\begin{tabularx}{\textwidth}{l p{3.0cm} X}
\toprule
\textbf{Plik} & \textbf{Warstwa} & \textbf{Opis wprowadzonych modyfikacji} \\
\midrule
\texttt{PackageItem.cs} & Core (Domain) & Nowa encja pozycji przesyłki (\texttt{Id}, \texttt{PackageId}, \texttt{Name}, \texttt{Price}, \texttt{ReturnPackageId}, \texttt{IsReturned}). \\
\addlinespace
\texttt{Package.cs} & Core (Domain) & Dodanie kolekcji \texttt{Items}, relacji \texttt{OriginalPackageId} / \texttt{ReturnDeliveries}, właściwości statusowych oraz metody \texttt{AddItem()}. \\
\addlinespace
\texttt{IPackageOrderService.cs} & Core (Domain) & Rozszerzenie kontraktu tworzenia paczki o opcjonalną listę deklarowanych produktów (\texttt{items}). \\
\addlinespace
\texttt{IPackageReturnService.cs} & Core (Domain) & Kontrakt operacji tworzenia przesyłek zwrotnych i zapytań o zwracalne towary. \\
\addlinespace
\texttt{DeliveryDbContext.cs} & Infrastructure & Dodanie \texttt{DbSet<PackageItem>} oraz mapowania relacji self-referencing w EF Core. \\
\addlinespace
\texttt{PackageRepository.cs} & Infrastructure & Eager loading pozycji \texttt{Items} oraz skrzynek paczkomatowych w \texttt{GetByIdAsync}. \\
\addlinespace
\texttt{PackageOrderService.cs} & Infrastructure & Zapisywanie zadeklarowanych produktów przy składaniu zamówienia. \\
\addlinespace
\texttt{PackageReturnService.cs} & Infrastructure & Orkiestracja zwrotu: walidacja uprawnień, blokada duplikatów, odwrócenie paczkomatów, alokacja przesyłki powrotnej. \\
\addlinespace
\texttt{Program.cs} (Cli \& Web) & Application / UI & Rejestracja \texttt{IPackageReturnService} w kontenerze Dependency Injection. \\
\addlinespace
\texttt{sequence-diagrams.puml} & Dokumentacja & Aktualizacja diagramu 01 o listę produktów oraz dodanie diagramu 06 dla zwrotu częściowego. \\
\addlinespace
\texttt{class-diagram.puml} & Dokumentacja & Dodanie encji \texttt{PackageItem}, relacji zwrotów \texttt{OriginalPackageId} i metod pomocniczych. \\
\bottomrule
\end{tabularx}
\end{table}

\section{Aktualizacja Diagramów Sekwencji i Modelu UML}

W katalogu \texttt{docs/} zaktualizowano dokumentację projektową w języku PlantUML:
\begin{enumerate}[leftmargin=*,itemsep=2pt]
    \item \textbf{Aktualizacja diagramu 01 (\texttt{01-order-and-payment}):}
    Diagram uwzględnia krok, w którym nadawca definiuje listę przesyłanych towarów (\texttt{items: [name, price]}). Metoda \texttt{CreatePackageAsync} przyjmuje kolekcję pozycji, które są atomowo zapisywane w bazie danych jako rekordy w tabeli \texttt{PackageItems} powiązane z nowo utworzoną przesyłką.
    \item \textbf{Nowy diagram 06 (\texttt{06-customer-return-items}):}
    Dodano kompletny diagram sekwencji dla procesu zwrotu (zarówno częściowego, jak i całkowitego):
    \begin{itemize}[itemsep=1pt]
        \item \textbf{Faza 1 (Pobranie zwracalnych pozycji):} Wywołanie \texttt{GetAvailableItemsForReturnAsync} weryfikuje tożsamość klienta (\texttt{RecipientId}) i zwraca wyłącznie produkty, dla których \texttt{IsReturned == false}.
        \item \textbf{Faza 2 (Utworzenie zwrotu):} Klient zaznacza zwracane pozycje. Metoda \texttt{CreateReturnDeliveryAsync} sprawdza status pierwotnej paczki (\texttt{Collected}), odwraca role (odbiorca staje się nadawcą zwrotu, a sklep odbiorcą) oraz paczkomaty, wywołuje standardowy \texttt{PackageOrderService} alokujący skrytki, powiązuje produkty kluczem \texttt{ReturnPackageId} i zleca asynchroniczne powiadomienie Hangfire.
        \item \textbf{Zasada ponownego użycia cyklu logistycznego:} Utworzona paczka zwrotna przechodzi przez ten sam istniejący pipeline kurierski (diagramy 02--05), co eliminuje redundancję kodu.
    \end{itemize}
    \item \textbf{Aktualizacja diagramu klas (\texttt{class-diagram.puml}):} Wprowadzono encję \texttt{PackageItem}, asocjację $1 \dots *$ z \texttt{Package} oraz rekurencyjną relację zwrotu (\texttt{OriginalPackage} / \texttt{ReturnDeliveries}).
\end{enumerate}

\newpage
\section{Zestawienie Scenariuszy Testowych}

Wszystkie testy wykonują się w pamięci RAM w ułamku sekundy (0.53 s), bez konieczności uruchamiania zewnętrznych baz danych:

\subsection{Zadanie 1: Testy procesu nadania i dostarczenia (\texttt{PackageLifecycleTests.cs})}
Zaimplementowano 20 scenariuszy (T01--T20) weryfikujących pełen cykl i przypadki skrajne:
\begin{itemize}[leftmargin=*,itemsep=1.5pt]
    \item \textbf{T01:} Pełny cykl Happy Path (\texttt{Ordered} $\rightarrow$ depozyt + wpłata $\rightarrow$ \texttt{Accepted} $\rightarrow$ \texttt{Preparing} $\rightarrow$ \texttt{Shipped} $\rightarrow$ \texttt{Delivered} $\rightarrow$ odbiór kodem PIN $\rightarrow$ \texttt{Collected}).
    \item \textbf{T02--T04:} Anulowanie na etapie \texttt{Ordered} i \texttt{Accepted} oraz blokada anulowania po wysyłce (\texttt{Shipped}).
    \item \textbf{T05:} Obsługa błędu bramki płatności i ponowienie próby zakończone sukcesem.
    \item \textbf{T06:} Pusta skrytka przy próbie odbioru przez kuriera --- natychmiastowe anulowanie zamówienia z winy nadawcy.
    \item \textbf{T07:} Przeterminowanie okna depozytowego (missed deposit job) i zwolnienie rezerwacji skrytki.
    \item \textbf{T08:} Zgubienie przesyłki w transporcie (upłynięcie 48h okna tranzytowego kuriera).
    \item \textbf{T09--T12:} Próba doręczenia w niewłaściwym statusie, błędny kod PIN, nieodebranie w 48h oraz próba odbioru przez nieuprawnioną osobę.
    \item \textbf{T13--T15:} Odrzucenie zamówienia przez operatora i blokada odrzucenia po wysyłce.
    \item \textbf{T16--T17:} Wymóg potwierdzenia depozytu i opłacenia paczki przed akceptacją przez operatora.
    \item \textbf{T18--T20:} Idempotencja potwierdzenia płatności, walidacja ocen kuriera (skala 1--5) i blokada przypisania kuriera w błędnym statusie.
\end{itemize}

\subsection{Zadanie 2: Moduł zwrotów częściowych (\texttt{PackageReturnTests.cs})}
Zaimplementowano 9 scenariuszy testowych weryfikujących nową funkcjonalność:
\begin{itemize}[leftmargin=*,itemsep=1.5pt]
    \item \textbf{R01:} \textbf{Zwrot częściowy (Happy Path)} --- zwrot 1 z 3 produktów tworzy nową paczkę zwrotną, zwracany produkt otrzymuje \texttt{ReturnPackageId}, a oryginalna paczka ma flagę \texttt{HasReturns == true}.
    \item \textbf{R02:} \textbf{Blokada przed odbiorem} --- zweryfikowano blokadę zwrotu dla statusów \texttt{Ordered}, \texttt{Accepted}, \texttt{Preparing}, \texttt{Shipped}.
    \item \textbf{R03:} \textbf{Weryfikacja uprawnień} --- zgłoszenie zwrotu przez użytkownika innego niż \texttt{RecipientId} rzuca wyjątek.
    \item \textbf{R04:} \textbf{Blokada ponownego zwrotu} --- próba zwrotu produktu, który już ma przypisany \texttt{ReturnPackageId}, jest odrzucana.
    \item \textbf{R05:} \textbf{Zwroty sekwencyjne w czasie} --- zwrot pierwszego produktu dzisiaj, a drugiego w późniejszym terminie w osobnej przesyłce zwrotnej.
    \item \textbf{R06:} \textbf{Automatyczne odwrócenie paczkomatów} --- paczka zwrotna jest nadawana w paczkomacie odbiorcy i kierowana do paczkomatu sklepu.
    \item \textbf{R07:} \textbf{Zwrot całkowity} --- odesłanie wszystkich produktów poprawnie ustawia właściwość \texttt{AreAllItemsReturned == true}.
    \item \textbf{R08:} \textbf{Tworzenie przesyłki z produktami} --- weryfikacja zapisu deklarowanej listy pozycji przez \texttt{CreatePackageAsync}.
    \item \textbf{R09:} \textbf{Metoda pomocnicza \texttt{AddItem()}} --- dodawanie produktów bezpośrednio do encji \texttt{Package}.
    \item \textbf{R10:} \textbf{Brak wolnej skrytki przy zwrocie} --- weryfikacja odmowy utworzenia zwrotu w sytuacji braku wolnej skrytki w paczkomacie zwrotu lub sklepu (towary pozostają niezwrócone, a wyjątek informuje klienta o konieczności wyboru innego paczkomatu/gabarytu).
\end{itemize}

\section{Podsumowanie Wyników Testów}

\begin{lstlisting}
Przebieg testu zakonczony powodzeniem.
Laczna liczba testow: 43
     Zakonczone pomyslnie: 43 (100% pass)
     Niepowodzenia: 0
     Pominiete: 0
Laczny czas: 0,85 sekundy
\end{lstlisting}

Wszystkie 43 testy (w tym 10 testów bazowych, 20 testów procesu nadania/dostarczenia T01--T20 oraz 13 testów modułu zwrotów R01--R10) kompilują się z zerową liczbą błędów i przechodzą pomyślnie.

\end{document}
